\section{Equality and the number of boundary components}
\label{q3:rigidity:setup}

Equality concerns the original metric and second fundamental form.
A numerical deformation can establish a sharp bound without producing
a limit that retains the original geometry.  We therefore use the
numerical comparison of Proposition~\ref{q3:intro:exterior-inequality}
to vary the original exterior, and then classify the stationary field
forced by equality.

There are two outermostness hypotheses below.  The finite-component
statement permits comparison cuts that retain some of the original
boundary components.  The connected statement needs only to exclude
cuts lying wholly in the open exterior.  We define these comparison
classes separately because their distinction is used in the proof.

\begin{definition}[Partial-coincidence outermostness]
\label{q3:rigidity:partial-outermostness}
Let the boundary of a one-ended exterior be
$S=\bigsqcup_{i=1}^{\ell}S_i$, with $1\leq\ell<\infty$ and each
$S_i$ connected. The boundary is outermost in the partial-coincidence
sense if no enclosing cut $\Gamma\ne S$ has all of the following
properties: every connected component of $\Gamma$ is either an entire
original component $S_i$ or is wholly contained in the open exterior;
at least one component lies in the open exterior; and
$\theta_+(\Gamma)\leq0$ everywhere for the normal toward the end.
\end{definition}

This condition includes comparisons that retain some original
components and replace others by interior components. It will be used
only after contact regularity has shown that the active minimizing
cuts have exactly this form. For a connected boundary, a separate
statement below retains the weaker condition excluding only wholly
interior enclosing cuts. The two comparison classes are not identified.

\begin{theorem}[Finite-component exterior equality rigidity]
\label{q3:intro:rigidity}
Let $(M,g,K)$ be a smooth connected orientable complete initial data
set without boundary, consisting of a compact core and finitely many
asymptotically flat ends. On each end assume
$g-\delta=O_2(r^{-q})$, $K=O_1(r^{-1-q})$, with $q>1/2$, and existing
finite ADM flux limits. Assume that $\mu$ and $|J|_g$ are integrable
and that $\mu\geq|J|_g$ everywhere.

Choose an end and a finite nonempty disjoint union
\[
S=\bigsqcup_{i=1}^{\ell}S_i,\qquad 1\leq\ell<\infty,
\]
of compact connected smooth embedded two-sided surfaces without boundary,
separating
that end from a complementary region. Let $\Mext$ be the closure of
the connected side toward the chosen end. Suppose that
$\partial\Mext=S$, that $\Mext$ has exactly this one end, and that,
for the normal pointing into $\Mext$:
\begin{enumerate}[label=\textup{(\roman*)}]
\item every $S_i$ is a future MOTS, so $H_{S_i}+\tr_{S_i}K=0$;
\item $S$ is outermost in the partial-coincidence sense of
Definition~\ref{q3:rigidity:partial-outermostness};
\item every enclosing cut in the closed exterior has area at least
$A=|S|_g=\sum_{i=1}^{\ell}|S_i|_g>0$.
\end{enumerate}
Let $(E,P)$ be the ADM vector of the chosen end. If
\begin{equation}
\label{q3:intro:equality}
E>|P|_\delta,\qquad
m=\sqrt{E^2-|P|_\delta^2}=\sqrt{\frac{A}{16\pi}},
\end{equation}
then $\ell=1$ and $S$ is diffeomorphic to a sphere. There is a proper
smooth spacelike embedding of the entire original $\Mext$ into the
maximally extended Schwarzschild spacetime of mass $m$, inducing $g$
and the given $K$ for a consistent future unit normal. The chosen end
approaches the corresponding spatial infinity. The image of $S$ is a
smooth cross-section of the corresponding future horizon or the
bifurcation sphere. The embedding and its future normal are smooth
up to $S$, and the image extends as a smooth spacelike hypersurface
through that boundary.
\end{theorem}

\begin{theorem}[Connected exterior equality rigidity]
\label{q3:intro:connected-rigidity}
Under the ambient, asymptotic, and exterior hypotheses of
Theorem~\ref{q3:intro:rigidity}, suppose that $S$ is connected.
Retain conditions \textup{(i)} and \textup{(iii)}, and replace
condition \textup{(ii)} by the following weaker outermostness condition:
there is no compact smooth embedded surface wholly contained in
$\operatorname{Int}\Mext$, possibly disconnected, separating $S$
from the distant end and having $\theta_+\le0$ everywhere for the
normal toward that end. If \eqref{q3:intro:equality} holds with
$A=|S|_g$, then $S$ is diffeomorphic to a sphere and all the
original-data Schwarzschild embedding conclusions of
Theorem~\ref{q3:intro:rigidity} hold.
\end{theorem}

\begin{corollary}[Strictness for disconnected boundary]
\label{q3:rigidity:finite-strictness}
Under the hypotheses of Theorem~\ref{q3:intro:rigidity}, retain
$E>|P|_\delta$ but do not assume its equality in mass and area.
If $\ell\geq2$, then
\[
\sqrt{E^2-|P|_\delta^2}>\sqrt{\frac{A}{16\pi}}.
\]
\end{corollary}

\begin{proof}
The original exterior is complete as a metric space with boundary,
as follows. An intrinsic Cauchy sequence is Cauchy in the complete
ambient manifold, and its ambient limit remains in the closed exterior.
A smooth interior or boundary half-ball chart then gives convergence
in the intrinsic distance. Its compact
boundary and unique end place it in
Definition~\ref{q3:intro:exterior-class}. Since $S$ is an enclosing cut
and every enclosing cut has area at least $A$, its enclosing infimum
is $A$. Proposition~\ref{q3:intro:exterior-inequality} gives the non-strict
inequality. Equality would imply $\ell=1$ by
Theorem~\ref{q3:intro:rigidity}, contradicting $\ell\geq2$.
\end{proof}

The proofs of Theorems~\ref{q3:intro:rigidity}
and~\ref{q3:intro:connected-rigidity} begin with the feasible strict
direction in Section~\ref{q3:strict:section}, continue through the
original-data variation and static base in
Sections~\ref{q3:variation:section} and~\ref{q3:static:section},
and conclude with global classification and horizon recovery in
Section~\ref{q3:global:section}.
Their outermostness assumptions enter only in localizing the area
multiplier in Lemma~\ref{q3:variation:area-support}, where the
connected case is treated separately. Each $S_i$ is orientable because it is two-sided in the
orientable manifold $M$, but no genus is prescribed. The proof first
constructs a causal stationary field on the original exterior and
concentrates its area multiplier on the entire boundary. A single
complete conformal-double argument then gives positive intrinsic
curvature and the same area $A$ on every component, forcing
$\ell=1$. Schwarzschild reconstruction follows with connectedness
established by the argument. The conclusion concerns only the chosen
exterior; it allows a nonconstant time graph and nonzero momentum in
the original asymptotic frame.

\paragraph{The proof path.}
First, the area derivative ranges over every minimizing enclosure.
Separation produces an area measure and causal constraint multipliers.
Normal-slice tests and the appropriate outermostness condition then
concentrate the area measure on the original boundary.  The resulting
stationary field can initially be null and can carry null dust.
We remove the twist without assuming regularity of that null set,
complete the static base, and apply a global spinor argument to its
conformal double.  This proves vacuum and recovers the original slice
through the horizon. After the global proof,
Section~\ref{q3:local-spinor:section} gives an independent local
conformal-family obstruction for two spherical components.
The strictness conclusion is qualitative; it supplies no positive
lower bound for the deficit depending only on the number of components.
